\chapter{Assembly Language Programming}

\section{From machine code to assembly}

Writing hex by hand is error-prone, and branch offsets must be recomputed every time code
is inserted. \textbf{Assembly language} is a textual, one-to-one representation of machine
code with two conveniences:
\begin{itemize}
  \item mnemonics and register names (\code{add \$3, \$1, \$2} instead of \code{0x00221820});
  \item \textbf{labels}: symbolic names for addresses, computed by the assembler.
\end{itemize}
Comments start with \code{;}. Each line has the form
\code{[labels:]* [instruction or .word]? [; comment]?}.

\section{Labels}

A label definition \code{name:} binds \code{name} to the address of the next instruction
(or \code{.word}) in the file. Several labels may precede the same instruction. A label may be
used:
\begin{itemize}
  \item in \code{beq}/\code{bne}, where the assembler computes the offset
        $(\mathit{label} - (\mathit{PC}_{\text{branch}} + 4))/4$;
  \item in \code{.word label}, which stores the label's \emph{absolute} address
        (used with \code{lis} + \code{jalr} to call procedures).
\end{itemize}

\begin{example}[Sum loop with labels]
The program from Chapter 2, now without hand-computed offsets:
\begin{lstlisting}[style=mips]
    lis $2
    .word 13
    lis $1
    .word 1
    add $3, $0, $0
top:
    add $3, $3, $2
    sub $2, $2, $1
    bne $2, $0, top
    jr $31
\end{lstlisting}
\end{example}

\section{Translating control structures}

\subsection{If-then-else}
For \code{if (a < b) S1 else S2} with $a$ in \code{\$1}, $b$ in \code{\$2}:
\begin{lstlisting}[style=mips]
    slt $5, $1, $2       ; $5 = a < b
    beq $5, $0, else     ; false -> else
    ; ... S1 ...
    beq $0, $0, endif    ; unconditional jump
else:
    ; ... S2 ...
endif:
\end{lstlisting}
\code{beq \$0, \$0, label} is the standard unconditional branch.
The resulting control flow:
\begin{center}
\begin{tikzpicture}[>=Stealth,node distance=0.5cm,font=\small,
  blk/.style={draw,rounded corners,minimum width=3.6cm,minimum height=0.6cm,align=center}]
  \node[blk] (t) {\code{slt \$5, \$1, \$2}};
  \node[blk,below=of t] (b) {\code{beq \$5, \$0, else}};
  \node[blk,below=of b,fill=blue!5] (s1) {S1};
  \node[blk,below=of s1] (j) {\code{beq \$0, \$0, endif}};
  \node[blk,below=of j,fill=blue!5] (s2) {S2};
  \node[blk,below=of s2,draw=none] (e) {\textit{next statement}};
  \node[font=\ttfamily\small,anchor=east] at ([xshift=-1.4cm]s2.west) {else:};
  \node[font=\ttfamily\small,anchor=east] at ([xshift=-1.4cm]e.west) {endif:};
  \draw[->] (t) -- (b);
  \draw[->] (b) -- node[right,font=\scriptsize]{$a<b$ (fall through)} (s1);
  \draw[->] (s1) -- (j);
  \draw[->] (b.west) -- ++(-1,0) |- node[pos=0.25,left,font=\scriptsize]{$a\ge b$} (s2.west);
  \draw[->] (j.east) -- ++(1,0) |- node[pos=0.25,right,font=\scriptsize]{always} (e.east);
  \draw[->] (s2) -- (e);
\end{tikzpicture}
\end{center}

\subsection{Comparisons}
Only \code{slt}/\code{sltu} exist; derive the rest:
\begin{itemize}
  \item $a > b$: \code{slt \$5, \$2, \$1} (swap operands).
  \item $a \ge b$: compute $a<b$ and invert. With \code{\$11} holding 1:
        \code{slt \$5,\$1,\$2; sub \$5,\$11,\$5}.
  \item $a \le b$: invert $b<a$.
  \item $a = b$, $a\neq b$: use \code{beq}/\code{bne} directly, or
        \code{slt}s: $(a<b) + (b<a)$ is $1$ iff $a\ne b$.
\end{itemize}

\subsection{Loops}
A \code{while (cond) body} becomes a test at the top and a jump back at the bottom:
\begin{lstlisting}[style=mips]
loop:
    ; compute cond into $5
    beq $5, $0, done
    ; ... body ...
    beq $0, $0, loop
done:
\end{lstlisting}

\section{Arrays}

With \code{mips.array}, \code{\$1} is the array address and \code{\$2} its length.
Element $i$ is at $\$1 + 4i$. Since \code{lw} takes only a constant offset, compute the
address explicitly.

\begin{example}[Maximum element of an array, result in \code{\$3}]
Run with \code{mips.array} (assumes length $\ge 1$):
\begin{lstlisting}[style=mips]
    lis $4
    .word 4            ; word size
    lis $11
    .word 1
    lw  $3, 0($1)      ; max = a[0]
    add $5, $1, $0     ; p = &a[0]
    add $6, $2, $0     ; n remaining
loop:
    lw  $7, 0($5)      ; x = *p
    slt $8, $3, $7     ; max < x ?
    beq $8, $0, skip
    add $3, $7, $0     ; max = x
skip:
    add $5, $5, $4     ; p++
    sub $6, $6, $11    ; n--
    bne $6, $0, loop
    jr $31
\end{lstlisting}
Pointer increment (``walk a pointer'') is cheaper than recomputing $\$1 + 4i$ each time.
To index directly: \code{mult \$i, \$4; mflo \$t; add \$t, \$t, \$1; lw \$x, 0(\$t)}.
\begin{center}
\begin{tikzpicture}[>=Stealth,cell/.style={draw,minimum width=1.6cm,minimum height=0.65cm,font=\small}]
  \node[cell] (c0) at (0,0) {$a[0]$};
  \node[cell,right=0pt of c0] (c1) {$a[1]$};
  \node[cell,right=0pt of c1,fill=blue!6] (c2) {$a[2]$};
  \node[cell,right=0pt of c2] (c3) {$\cdots$};
  \node[cell,right=0pt of c3] (c4) {$a[n-1]$};
  \foreach \c/\t in {c0/\$1,c1/\$1+4,c2/\$1+8,c4/\$1+4(n-1)} \node[below,font=\scriptsize] at (\c.south) {\code{\t}};
  \draw[<-,thick,blue!60!black] (c2.north) -- ++(0,0.8) node[above,font=\small]{\code{\$5} ($p$)};
  \draw[->,gray] (c1.north) to[bend left=50] node[above,font=\scriptsize]{$p \mathrel{+}= 4$} (c2.north west);
  \draw[<-] (c0.west) -- ++(-0.8,0) node[left,font=\small]{\code{\$1}};
\end{tikzpicture}
\end{center}
\end{example}

\section{The stack}

Registers are a shared global resource; procedures need private storage. The loader
initialises \code{\$30} to the end of memory and we use it as a \textbf{stack pointer}. The
stack grows \emph{downward}, and \code{\$30} points to the last-pushed word (the top).

\begin{lstlisting}[style=mips]
; push $5
sw  $5, -4($30)
sub $30, $30, $4      ; $4 = 4
; pop into $5
add $30, $30, $4
lw  $5, -4($30)
\end{lstlisting}

\begin{center}
\begin{tikzpicture}[cell/.style={draw,minimum width=3cm,minimum height=0.6cm,font=\small}]
  \node[cell] (a) at (0,0) {caller's data};
  \node[cell,below=0pt of a] (b) {saved \code{\$31}};
  \node[cell,below=0pt of b] (c) {saved \code{\$5}};
  \node[cell,below=0pt of c,fill=gray!10] (d) {free};
  \node[cell,below=0pt of d,fill=gray!10] (e) {free};
  \draw[Stealth-] (c.east) -- ++(1.2,0) node[right]{\code{\$30}};
  \node[left] at (a.north west) {\small high addresses};
  \node[left] at (e.south west) {\small low addresses};
  \draw[-Stealth] ($(a.west)+(-2.6,0)$) -- ($(e.west)+(-2.6,0)$) node[midway,left,font=\small]{grows};
\end{tikzpicture}
\end{center}

\section{Procedures}

A procedure is a labelled block of code that ends in \code{jr \$31}. To call it:
\begin{lstlisting}[style=mips]
    lis $5
    .word proc       ; absolute address of proc
    jalr $5          ; $31 = return address, jump to proc
\end{lstlisting}

\paragraph{The \code{\$31} problem.} \code{jalr} overwrites \code{\$31}, which our own
program needs to return to the loader. So every caller (including \code{main}) must save
\code{\$31} on the stack before \code{jalr} and restore it after:
\begin{lstlisting}[style=mips]
    sw   $31, -4($30)    ; push $31
    lis  $31
    .word 4
    sub  $30, $30, $31
    lis  $5
    .word proc
    jalr $5
    lis  $31
    .word 4
    add  $30, $30, $31   ; pop $31
    lw   $31, -4($30)
\end{lstlisting}
(Here \code{\$31} itself is reused as a temporary holding 4, since its real value is safely
on the stack.)

\subsection{Calling conventions}
A \textbf{calling convention} is the contract between caller and callee. The convention we
use in CS 241 assignments:
\begin{itemize}
  \item Arguments are passed in registers (documented per procedure) or on the stack.
  \item The return value is placed in \code{\$3}.
  \item \textbf{Callee-save}: the procedure must restore every register it modifies, except
        its documented output(s). It pushes them at entry and pops them before \code{jr \$31}.
  \item The caller saves \code{\$31} around \code{jalr}.
\end{itemize}
The alternative, \emph{caller-save}, makes the caller save any register it cares about
before calling. Real ABIs mix the two (some registers callee-saved, some caller-saved).

\begin{important}[Document your procedures]
For each procedure write a header comment: what it does, which registers hold inputs,
which hold outputs, and which (if any) are clobbered. Bugs in assembly are almost always
convention violations: a forgotten restore, a stack pointer that is off by 4, or \code{\$31}
not saved.
\end{important}

\begin{example}[Procedure: sum of 1..\code{\$1} in \code{\$3}]
A callee-save procedure: it restores every register it uses except its output.
\begin{lstlisting}[style=mips]
; sum1toN: $3 = 1 + 2 + ... + $1 ($1 >= 1)
; registers modified: $3 (output) -- $1, $2 preserved
sum1toN:
    sw  $1, -4($30)      ; save $1
    sw  $2, -8($30)      ; save $2
    lis $2
    .word 8
    sub $30, $30, $2     ; reserve 2 words
    lis $2
    .word 1
    add $3, $0, $0
sloop:
    add $3, $3, $1
    sub $1, $1, $2
    bne $1, $0, sloop
    lis $2
    .word 8
    add $30, $30, $2     ; release 2 words
    lw  $1, -4($30)
    lw  $2, -8($30)
    jr  $31
\end{lstlisting}
\end{example}

\section{Stack frames and recursion}

Recursion works automatically once every procedure saves \code{\$31} and its registers on
the stack: each activation gets its own region, its \textbf{stack frame} (activation record).
A frame typically contains saved registers, the saved return address, and local variables.

\begin{example}[Recursive factorial]
Each activation pushes a three-word frame holding \code{\$1}, \code{\$2} and \code{\$31}.
\begin{lstlisting}[style=mips]
; fact: $3 = $1!   (uses $1 as input; preserves all but $3)
fact:
    sw  $1,  -4($30)
    sw  $2,  -8($30)
    sw  $31, -12($30)
    lis $2
    .word 12
    sub $30, $30, $2     ; push frame of 3 words
    lis $3
    .word 1              ; base result
    lis $2
    .word 1
    slt $2, $2, $1       ; $2 = (1 < n)
    beq $2, $0, fdone    ; n <= 1 -> return 1
    lis $2
    .word 1
    sub $1, $1, $2       ; n - 1
    lis $2
    .word fact
    jalr $2              ; $3 = (n-1)!
    lw  $1, 8($30)       ; reload original n from frame
    mult $3, $1
    mflo $3              ; $3 = n * (n-1)!
fdone:
    lis $2
    .word 12
    add $30, $30, $2     ; pop frame
    lw  $1,  -4($30)
    lw  $2,  -8($30)
    lw  $31, -12($30)
    jr  $31
\end{lstlisting}
After pushing, the saved $\$1$ is at offset 8 from the new \code{\$30} (frame is
$\$30{+}0$: \code{\$31}, $\$30{+}4$: \code{\$2}, $\$30{+}8$: \code{\$1}).
The stack at the deepest point of computing $3!$ (called from \code{main}) holds three frames:
\begin{center}
\begin{tikzpicture}[>=Stealth,cell/.style={draw,minimum width=4.4cm,minimum height=0.5cm,font=\small}]
  \node[cell,fill=gray!10] (top) at (0,0) {caller's data};
  \node[cell,below=0pt of top] (a1) {\code{\$1} $=3$};
  \node[cell,below=0pt of a1] (a2) {\code{\$2} (caller's value)};
  \node[cell,below=0pt of a2] (a3) {\code{\$31}: return into \code{main}};
  \node[cell,below=0pt of a3,fill=blue!4] (b1) {\code{\$1} $=2$};
  \node[cell,below=0pt of b1,fill=blue!4] (b2) {\code{\$2} $=$ address of \code{fact}};
  \node[cell,below=0pt of b2,fill=blue!4] (b3) {\code{\$31}: return into \code{fact}};
  \node[cell,below=0pt of b3] (c1) {\code{\$1} $=1$};
  \node[cell,below=0pt of c1] (c2) {\code{\$2} $=$ address of \code{fact}};
  \node[cell,below=0pt of c2] (c3) {\code{\$31}: return into \code{fact}};
  \draw[decorate,decoration={brace,amplitude=5pt}] ([xshift=2pt]a1.north east) -- ([xshift=2pt]a3.south east) node[midway,right=6pt,font=\small]{frame of \code{fact(3)}};
  \draw[decorate,decoration={brace,amplitude=5pt}] ([xshift=2pt]b1.north east) -- ([xshift=2pt]b3.south east) node[midway,right=6pt,font=\small]{frame of \code{fact(2)}};
  \draw[decorate,decoration={brace,amplitude=5pt}] ([xshift=2pt]c1.north east) -- ([xshift=2pt]c3.south east) node[midway,right=6pt,font=\small]{frame of \code{fact(1)}};
  \draw[<-] (c3.west) -- ++(-1,0) node[left]{\code{\$30}};
\end{tikzpicture}
\end{center}
Each frame is popped as its activation returns, and \code{\$3} carries $1$, then $2\cdot1$, then $3\cdot2$ back up.
\end{example}

\paragraph{Recursion depth.} Each call uses a fixed amount of stack, so depth $d$ uses
$\Theta(d)$ stack words. Infinite recursion eventually runs \code{\$30} into the program or
heap and crashes.

\section{Input and output}

\begin{example}[Echo standard input to standard output]
Read bytes until EOF ($-1$), writing each one back out.
\begin{lstlisting}[style=mips]
    lis $1
    .word 0xffff0004     ; input
    lis $2
    .word 0xffff000c     ; output
    lis $4
    .word -1             ; EOF marker
loop:
    lw  $3, 0($1)        ; read a byte
    beq $3, $4, end
    sw  $3, 0($2)        ; write it
    beq $0, $0, loop
end:
    jr $31
\end{lstlisting}
\end{example}

\paragraph{Printing a number.} To print a signed integer: print \code{-} and negate if
negative; then repeatedly \code{div} by 10, pushing remainders (\code{mfhi}) on the stack
until the quotient is 0; finally pop each digit, add 48, and store it to \code{0xffff000c}.
Pitfalls: the number 0 must print one digit, and $-2^{31}$ cannot be negated (handle it via
unsigned division or by working with negative remainders).

\section{Summary}
\begin{itemize}
  \item Assembly = machine code with mnemonics, comments and labels; the assembler computes
        branch offsets and label addresses.
  \item \code{if}/\code{while} translate to \code{slt} + \code{beq}/\code{bne};
        \code{beq \$0,\$0,L} is an unconditional jump.
  \item The stack grows down from \code{\$30}; push = store then decrement, pop = increment then load.
  \item Call with \code{lis}/\code{.word}/\code{jalr}; save \code{\$31} first. Return with \code{jr \$31}.
  \item Calling convention: callee restores everything except outputs; results in \code{\$3}.
  \item Recursion needs nothing extra: each activation's state lives in its own stack frame.
\end{itemize}
